\subsection{Explicit $\Ethree$ operations for
$V_k(\mathfrak{sl}_2)$}
\label{subsec:e3-explicit}

The abstract $\Ethree$ structure of
Proposition~\ref{prop:e3-structure} acquires concrete content
for $\fg = \mathfrak{sl}_2$: the five-dimensional derived
chiral centre $\CC \oplus \mathfrak{sl}_2[-1] \oplus \CC[-2]$
carries an explicit cup product, Gerstenhaber bracket,
$\Pthree$ bracket, and BV operator, all computable from the
OPE.

\medskip
\noindent
\textbf{The underlying graded vector space.}
By Proposition~\ref{prop:e3-structure} and the chiral
Hochschild computation of
Example~\ref{ex:km}, the derived chiral centre of
$V_k(\mathfrak{sl}_2)$ at generic level
$k \neq -2$ is the graded vector space
\begin{equation}\label{eq:e3-graded-space}
  Z^{\mathrm{der}}_{\mathrm{ch}}(V_k(\mathfrak{sl}_2))
  \;=\;
  \underbrace{\CC}_{=\, \HH^0}
  \;\oplus\;
  \underbrace{\mathfrak{sl}_2[-1]}_{=\, \HH^1}
  \;\oplus\;
  \underbrace{\CC[-2]}_{=\, \HH^2},
\end{equation}
with total dimension $1 + 3 + 1 = 5$. Write $\mathbf{1}$
for the generator of $\HH^0 = \CC$ (the vacuum),
$\{e, f, h\}$ for the standard basis of
$\HH^1 = \mathfrak{sl}_2$
(outer derivations: infinitesimal level-preserving
deformations parametrised by $\fg$ itself), and $\eta$
for the generator of $\HH^2 = \CC$
(the central charge deformation cocycle). The
$\mathfrak{sl}_2$-module structure is: $\HH^0$ trivial,
$\HH^1$ adjoint, $\HH^2$ trivial.

\medskip
\noindent
\textbf{Conventions.}
The Killing form normalisation is
$(e, f) = (f, e) = 1$, $(h, h) = 2$, all other pairings zero.
Structure constants: $[e, f] = h$, $[h, e] = 2e$,
$[h, f] = -2f$. The inverse Casimir is
$\Omega = e \otimes f + f \otimes e
+ \tfrac{1}{2}\,h \otimes h$, and $h^\vee = 2$.
The KZ coupling constant is
$h_{\mathrm{KZ}} = 1/(k + h^\vee) = 1/(k + 2)$.

\begin{remark}[Convention for $\hbar$]
\label{rem:hbar-convention}
Throughout this paper, $\hbar$ denotes the
\emph{KZ coupling constant}
\begin{equation}\label{eq:hbar-convention}
  \hbar \;=\; h_{\mathrm{KZ}}
  \;=\; \frac{1}{k + h^\vee},
\end{equation}
the normalisation in the KZ connection
$\nabla = d + \hbar \sum \Omega_{ij}\,dz_{ij}$.
It vanishes as $k \to \infty$ (weak-coupling limit) and
diverges at the critical level $k = -h^\vee$.
In the $\mathfrak{sl}_2$ computations, $h^\vee = 2$,
so $\hbar = 1/(k+2)$.
The \emph{departure from critical level} is the inverse
quantity $k + h^\vee = 1/\hbar$; we always write this
explicitly as $k + h^\vee$ rather than introducing a
separate symbol. Critical level corresponds to
$k + h^\vee = 0$ (equivalently, $\hbar \to \infty$);
weak coupling corresponds to $k + h^\vee \to \infty$
(equivalently, $\hbar \to 0$).
For the quantum group parameter: $q = e^{2\pi i \hbar}$.
\end{remark}

An $\Ethree$-algebra carries three levels of
operations:
\begin{itemize}
\item Level $0$ ($\Einf$): a graded-commutative product
  $\mu\colon \HH^p \otimes \HH^q \to \HH^{p+q}$.
\item Level $1$ ($\Etwo$ beyond $\Eone$): the
  Gerstenhaber bracket
  $[-,-]\colon \HH^p \otimes \HH^q \to \HH^{p+q-1}$
  of degree~$-1$, derived from the fundamental class of
  $S^1 = \Conf_2(\RR^2)/\RR^+$.
\item Level $2$ ($\Ethree$ beyond $\Etwo$): the $\Pthree$
  bracket
  $\{-,-\}\colon \HH^p \otimes \HH^q \to \HH^{p+q-2}$
  of degree~$-2$, derived from the fundamental class of
  $S^2 = \Conf_2(\RR^3)/\RR^+$.
\end{itemize}
We compute each in turn.

\begin{proposition}[Explicit $\Ethree$ operations on
$Z^{\mathrm{der}}_{\mathrm{ch}}(V_k(\mathfrak{sl}_2))$]
\label{prop:e3-explicit-sl2}
Let $k \neq -2$.
The $\Ethree$-algebra structure on
$Z^{\mathrm{der}}_{\mathrm{ch}}(V_k(\mathfrak{sl}_2))$
is determined by the following data.

\smallskip
\noindent
\textup{(I) Cup product} ($\Einf$ level).
\begin{enumerate}[label=\textup{(\alph*)}]
\item $\mathbf{1}$ is the unit:
  $\mathbf{1} \cdot f = f \cdot \mathbf{1} = f$ for all $f$.
\item
  $\mu(X, Y) = 0$ for all $X, Y \in \HH^1 = \mathfrak{sl}_2$.
\item
  $\mu(X, \eta) = 0$ for all $X \in \HH^1$, and
  $\eta \cdot \eta = 0$.
\end{enumerate}

\smallskip
\noindent
\textup{(II) Gerstenhaber bracket} ($\Etwo$ level,
degree~$-1$).
\begin{gather*}
  [\mathbf{1}, -] = 0, \qquad
  [X, Y] = 0 \quad \text{for all }
  X, Y \in \HH^1 = \mathfrak{sl}_2, \\
  [X, \eta] = 0 \quad \text{for all } X \in \HH^1,
  \qquad [\eta, \eta] = 0.
\end{gather*}

\smallskip
\noindent
\textup{(III) The $\Pthree$ bracket} ($\Ethree$ level,
degree~$-2$).
\begin{align}
  \label{eq:e3-p3-killing}
  \{X, Y\}
  &= h_{\mathrm{KZ}} \cdot (X, Y)
  = \frac{1}{k + h^\vee}\,(X, Y)
  \quad \text{for } X, Y \in \HH^1,
  \\
  \label{eq:e3-p3-derivation}
  \{X, \eta\}
  &= 0
  \quad \text{for } X \in \HH^1,
  \\
  \label{eq:e3-p3-unit}
  \{\mathbf{1}, -\}
  &= 0,
  \\
  \label{eq:e3-p3-top}
  \{\eta, \eta\}
  &= 0,
\end{align}
where $(X, Y)$ denotes the normalised Killing form
and $h_{\mathrm{KZ}} = 1/(k + h^\vee) = 1/(k + 2)$ is
the KZ coupling constant
\textup{(}Remark~\textup{\ref{rem:hbar-convention}}\textup{)}.
\end{proposition}

\begin{proof}
We establish each level by $\mathfrak{sl}_2$-equivariance,
degree counting, and the BV relation.

\textit{Cup product.}
The cup product $\mu\colon \HH^p \otimes \HH^q
\to \HH^{p+q}$ is graded-commutative:
$\mu(a, b) = (-1)^{pq}\,\mu(b, a)$.
For $X, Y \in \HH^1$, this gives
$\mu(X, Y) = -\mu(Y, X)$: the restriction to
$\HH^1 \otimes \HH^1 \to \HH^2 = \CC$ is an
$\mathfrak{sl}_2$-invariant antisymmetric bilinear form
on the adjoint representation valued in the trivial
representation.
Since
$\Lambda^2(\mathfrak{sl}_2)^{\mathfrak{sl}_2}
\cong \Hom_{\mathfrak{sl}_2}(\Lambda^2(\mathrm{ad}), \CC)
= 0$
(the exterior square of the adjoint decomposes as
$\Lambda^2(\mathrm{ad}) \cong \mathrm{ad}$,
which has no trivial summand), the cup product
$\mu(X, Y) = 0$ for all $X, Y \in \HH^1$.

The products $\mu(\HH^1, \HH^2) \subset \HH^3 = 0$ and
$\mu(\HH^2, \HH^2) \subset \HH^4 = 0$ vanish for degree
reasons.

\textit{BV operator.}
The framed $\Etwo$ structure on $\BarSig(V_k(\mathfrak{sl}_2))$
(from genus-$1$ sewing) equips the derived centre with a BV
operator $\Delta\colon \HH^n \to \HH^{n-1}$ satisfying
$\Delta^2 = 0$ and the BV relation
\begin{equation}\label{eq:bv-rel-sl2}
  [f, g]
  = \Delta(f \cdot g) - (\Delta f) \cdot g
  - (-1)^{|f|}\,f \cdot (\Delta g).
\end{equation}

The component $\Delta\colon \HH^1 \to \HH^0$ is an
$\mathfrak{sl}_2$-equivariant linear map from the adjoint
representation to the trivial representation. Since the
adjoint is irreducible and non-trivial,
$\Delta|_{\HH^1} = 0$.

The component $\Delta\colon \HH^2 \to \HH^1$ is an
$\mathfrak{sl}_2$-equivariant linear map from the trivial
to the adjoint. Again, $\Delta|_{\HH^2} = 0$.

All other components vanish trivially
($\Delta|_{\HH^0}$ lands in $\HH^{-1} = 0$;
$\Delta|_{\HH^n} = 0$ for $n \geq 3$). Therefore
$\Delta = 0$ on the entire derived centre.

\textit{Gerstenhaber bracket.}
With $\Delta = 0$ and $\mu|_{\HH^1 \otimes \HH^1} = 0$,
the BV relation~\eqref{eq:bv-rel-sl2} gives:
for $X, Y \in \HH^1$,
\[
  [X, Y]
  = \Delta(\mu(X, Y)) - (\Delta X) \cdot Y
  + X \cdot (\Delta Y)
  = \Delta(0) - 0 + 0 = 0.
\]
For $X \in \HH^1$, $\eta \in \HH^2$:
$[X, \eta] = \Delta(\mu(X, \eta)) - (\Delta X) \cdot \eta
- (-1)^1\,X \cdot (\Delta \eta)
= 0 - 0 + 0 = 0$
(since $\mu(X, \eta) \in \HH^3 = 0$,
$\Delta X = 0$, $\Delta\eta = 0$).
Similarly $[\eta, \eta] \in \HH^3 = 0$.

The vanishing of the Gerstenhaber bracket on the derived
centre is a structural consequence of the
$\mathfrak{sl}_2$-equivariance constraints: the adjoint
representation provides no equivariant maps to or from the
trivial representation. This is specific to the
\emph{derived centre} $\HH^*(\BarSig, \BarSig)$; the
Gerstenhaber bracket on the full Hochschild cochain complex
$C^*(\BarSig, \BarSig)$ is non-trivial (it carries the
commutator of derivations at the cochain level).

\textit{$\Pthree$ bracket.}
The degree-$(-2)$ bracket
$\{-,-\}\colon \HH^p \otimes \HH^q \to \HH^{p+q-2}$
satisfies the graded symmetry
\begin{equation}\label{eq:p3-symmetry}
  \{a, b\}
  = -(-1)^{(|a|-2)(|b|-2)}\,\{b, a\}
\end{equation}
of a degree-$(-2)$ Lie bracket, together with the Leibniz rule
$\{a, b \cdot c\}
= \{a, b\} \cdot c + (-1)^{(|a|-2)|b|}\,b \cdot \{a, c\}$.

\textit{Step 1: $\{\mathbf{1}, -\} = 0$.}
The Leibniz rule applied to
$\mathbf{1} = \mathbf{1} \cdot \mathbf{1}$ gives
$\{\mathbf{1}, f\} = 2\,\{\mathbf{1}, f\}$
for all $f$; hence $\{\mathbf{1}, f\} = 0$.

\textit{Step 2: $\{X, Y\} = h_{\mathrm{KZ}}\,(X, Y)$.}
For $X, Y \in \HH^1$,
$\{X, Y\} \in \HH^0 = \CC$.
The symmetry~\eqref{eq:p3-symmetry} at
$|X| = |Y| = 1$ gives
$\{X, Y\} = -(-1)^{(-1)(-1)}\,\{Y, X\}
= -(-1)^1\,\{Y, X\}
= \{Y, X\}$:
the bracket is \emph{symmetric} on $\HH^1 \otimes \HH^1$.
By $\mathfrak{sl}_2$-equivariance, $\{X, Y\}$ is an
invariant symmetric bilinear form on the adjoint
representation valued in $\CC$. The space of such forms is
one-dimensional, spanned by the Killing form $(X, Y)$.
Hence $\{X, Y\} = \alpha \cdot (X, Y)$ for a scalar
$\alpha$ depending on $k$.

The scalar $\alpha$ is determined by the $r$-matrix
residue, independently of any comparison with the CFG
construction. The $\Pthree$ bracket on the derived
chiral centre $Z^{\mathrm{der}}_{\mathrm{ch}}(V_k(\fg))
= \HH^*(\BarSig(V_k(\fg)), \BarSig(V_k(\fg)))$ is the
$\Etwo$-Hochschild cohomology operation derived from
the fundamental class of $S^2 \subset \Conf_2(\RR^3)$.
At degree~$2$, this operation is the residue of the
classical $r$-matrix paired with the fundamental class
of $S^1 \subset \Conf_2(\CC)$. In the KZ normalisation
(see~\eqref{eq:kz-sl2-degree2} and the conventions
of~\S\ref{subsec:sl2-chiral-e3}):
% AP126/AP148: KZ convention r(z) = Omega/((k+h^v)z).
% At k=0, r = Omega/(h^v z) != 0 (non-abelian; correct for KZ).
% At k=-h^v, r diverges (Sugawara singularity).
\begin{equation}\label{eq:p3-from-r-matrix-residue}
  r(z) = \frac{\Omega}{(k + h^\vee)\,z},
  \qquad
  \operatorname{Res}_{z=0}\bigl[r(z)\bigr]
  = \frac{\Omega}{k + h^\vee}.
\end{equation}
The residue acts on $X \otimes Y \in \fg \otimes \fg$
by contraction with the inverse Casimir:
$\Omega(X \otimes Y) = (X, Y)$ (the normalised Killing
form). The $\Pthree$ bracket on $\HH^1 \cong \fg$ is
therefore
\begin{equation}\label{eq:p3-bracket-normalisation}
  \{X, Y\}
  = \operatorname{Res}_{z=0}
  \bigl[r(z)(X \otimes Y)\bigr]
  = \frac{(X, Y)}{k + h^\vee}
  = h_{\mathrm{KZ}} \cdot (X, Y),
\end{equation}
giving $\alpha = h_{\mathrm{KZ}} = 1/(k + h^\vee)$.
For $\mathfrak{sl}_2$, $h^\vee = 2$ and
$\alpha = 1/(k + 2)$.

Verification: at $k \to \infty$,
$h_{\mathrm{KZ}} \to 0$ and the bracket vanishes (the
weak-coupling limit is the abelian $\Ethree$). At the
critical level $k = -h^\vee$, $h_{\mathrm{KZ}}$ diverges
(the bracket degenerates, reflecting the Feigin--Frenkel
singularity). Both limits are consistent with
the behaviour of the KZ connection.

\textit{Step 3: $\{X, \eta\} = 0$.}
The bracket $\{X, \eta\} \in \HH^1 = \mathfrak{sl}_2$
for $X \in \HH^1$, $\eta \in \HH^2$ is an
$\mathfrak{sl}_2$-equivariant map
$\mathrm{ad} \otimes \CC \to \mathrm{ad}$,
hence $\{X, \eta\} = \beta\,X$ for some scalar $\beta$.
The $\Pthree$ Jacobi identity applied to
$X, Y \in \HH^1$ and $\eta \in \HH^2$,
combined with
$\{X, Y\} = h_{\mathrm{KZ}}\,(X, Y) \in \HH^0$, gives
$\{\{X, \eta\}, \eta\} = 0$ (since
$\{\{X, \eta\}, \eta\} \in \HH^{-1} = 0$) and
$\{\eta, \{X, \eta\}\} \in \HH^{-1} = 0$,
while $\{X, \{\eta, \eta\}\} = \{X, 0\} = 0$.
The Jacobi identity is thus satisfied for all $\beta$;
the constraint comes instead from the self-consistency
of the $\Ethree$ deformation.
At the cohomological level, one applies the iterated Jacobi
identity to $X, Y \in \HH^1$ and $\eta$:
\begin{align*}
  \{\{X, Y\}, \eta\}
  &= \{X, \{Y, \eta\}\}
  - (-1)^{(|X|-2)(|Y|-2)}\,\{Y, \{X, \eta\}\} \\
  &= \{X, \beta Y\} - (-1)^1\,\{Y, \beta X\} \\
  &= \beta\,h_{\mathrm{KZ}}\,(X, Y)
  + \beta\,h_{\mathrm{KZ}}\,(Y, X) \\
  &= 2\beta\,h_{\mathrm{KZ}}\,(X, Y).
\end{align*}
The left side:
$\{\{X, Y\}, \eta\}
= \{h_{\mathrm{KZ}}(X, Y)\,\mathbf{1}, \eta\}
= h_{\mathrm{KZ}}(X, Y)\,\{\mathbf{1}, \eta\} = 0$.
Hence $2\beta\,h_{\mathrm{KZ}}\,(X, Y) = 0$ for all
$X, Y$. Since
$(e, f) = 1 \neq 0$ and $h_{\mathrm{KZ}} \neq 0$ at
generic $k$, we conclude $\beta = 0$.

\textit{Step 4: $\{\eta, \eta\} = 0$.}
The bracket $\{\eta, \eta\} \in \HH^2 = \CC$ and
the symmetry~\eqref{eq:p3-symmetry} at
$|\eta| = 2$ gives
$\{\eta, \eta\} = -(-1)^{0 \cdot 0}\,\{\eta, \eta\}
= -\{\eta, \eta\}$, hence $\{\eta, \eta\} = 0$.
\end{proof}

\begin{remark}[The Killing form as the $\Ethree$ shadow
of the Cartan $3$-form]
\label{rem:cartan-3-form}
The bracket~\eqref{eq:e3-p3-killing} has a transparent
topological origin. The $\Ethree$-deformation space is
$(k+h^\vee)\,H^3(\mathfrak{sl}_2)[[k+h^\vee]] \cong (k+h^\vee)\,\CC[[k+h^\vee]]$,
one-dimensional at each order in $(k+h^\vee)$
(Theorem~\ref{thm:cfg}(iv)). The generator of
$H^3(\mathfrak{sl}_2)$ is the Cartan $3$-form
\begin{equation}\label{eq:cartan-3form}
  \omega_3(X, Y, Z)
  = (X, [Y, Z]),
\end{equation}
the unique (up to scalar) $\mathfrak{sl}_2$-invariant element
of $\Lambda^3(\mathfrak{sl}_2^*)$. The degree-$(-2)$ bracket
$\{X, Y\}$ is the $\Ethree$-operation induced by this
deformation class scaled by
$\hbar = 1/(k + h^\vee)$:
\[
  \{X, Y\}
  = \hbar \cdot (X, Y).
\]
The mechanism is \emph{not} literal contraction of the
antisymmetric $3$-form $\omega_3$: the $\Pthree$ bracket
is symmetric on $\HH^1$ (by the graded
symmetry~\eqref{eq:p3-symmetry} at $|X| = |Y| = 1$),
while any contraction of $\omega_3$ in two slots is
antisymmetric. The connection is indirect:
the one-dimensional deformation space
$H^3(\mathfrak{sl}_2)$, generated by $\omega_3$,
parametrises $\Ethree$-deformations of $C^*(\fg)$
(Theorem~\ref{thm:cfg}(iv));
on cohomology, $\mathfrak{sl}_2$-equivariance forces
the induced $\Pthree$ bracket on
$\HH^1 \cong \mathfrak{sl}_2$ to be proportional to
the Killing form $(X, Y)$, the unique invariant
symmetric bilinear form. The coefficient
$h_{\mathrm{KZ}} = 1/(k + h^\vee)$ is fixed by the
$r$-matrix residue
(equation~\eqref{eq:p3-from-r-matrix-residue}), not by
comparison with CFG.
The $\Pthree$ bracket on $Z^{\mathrm{der}}_{\mathrm{ch}}$
is the binary shadow of the ternary Cartan deformation class,
with the Killing form appearing through equivariance rather
than through direct contraction.
\end{remark}

% ----------------------------------------------------------------
\subsection{Quantum corrections from the vertex $R$-matrix:
$\Ethree$ operations for the EK quantum VOA}
\label{subsec:e3-ek-quantum}

The classical $V_k(\mathfrak{sl}_2)$ has trivial braiding
$S(z) = \id$; the Etingof--Kazhdan quantum VOA
$V_{\mathrm{EK}}$ of Example~\ref{ex:ek-qvoa} has
$S(z) = \id + h_{\mathrm{KZ}}\,\Omega/z + O(z^{-2})$. Does
the non-trivial $R$-matrix change the $\Ethree$-algebra
structure on the derived chiral centre? On cohomology, no:
the $\Pthree$ bracket coincides with the classical answer
(Proposition~\ref{prop:e3-explicit-sl2}). On cochains, yes:
the quantum corrections encode the braided monoidal structure
of $\mathrm{Rep}(Y_\hbar(\mathfrak{sl}_2))$.

\medskip
\noindent
\textbf{The underlying graded vector space.}
The EK quantum VOA is a flat formal deformation of
$V_k(\mathfrak{sl}_2)$ over $\CC[[h_{\mathrm{KZ}}]]$, with
$h_{\mathrm{KZ}} = 1/(k+2)$.
Since the chiral Hochschild cohomology is computed by
a bounded complex of finite-dimensional
$\mathfrak{sl}_2$-modules, the deformation does not
change the underlying graded vector space:
\begin{equation}\label{eq:ek-e3-graded}
  Z^{\mathrm{der}}_{\mathrm{ch}}(V_{\mathrm{EK}})
  \;=\;
  \CC[[h_{\mathrm{KZ}}]]
  \;\oplus\;
  \mathfrak{sl}_2[[h_{\mathrm{KZ}}]][-1]
  \;\oplus\;
  \CC[[h_{\mathrm{KZ}}]][-2],
\end{equation}
with total rank $5$ over $\CC[[h_{\mathrm{KZ}}]]$.
The generators
$\mathbf{1} \in \HH^0$, $\{e, f, h\} \in \HH^1$, and
$\eta \in \HH^2$ are defined as in
Proposition~\ref{prop:e3-explicit-sl2}, now taken
over the base ring $\CC[[h_{\mathrm{KZ}}]]$.

\medskip
\noindent
\textbf{Equivariance under the quantum $R$-matrix.}
The Yang $R$-matrix $R(u) = u\,I + \hbar\,P$ on
$V \otimes V$ is $\mathfrak{sl}_2$-equivariant:
$[R(u),\, \rho(X) \otimes I + I \otimes \rho(X)] = 0$
for all $X \in \mathfrak{sl}_2$.
The full vertex $R$-matrix
$S(z)$ on $V_{\mathrm{EK}}^{\otimes 2}$
(equation~\eqref{eq:ek-vertex-rmatrix}) is likewise
$\mathfrak{sl}_2$-equivariant, since it is constructed from
$\mathcal{R}(z)$, the spectral universal $R$-matrix of
the Yangian $Y_\hbar(\mathfrak{sl}_2)$, which
commutes with the $\mathfrak{sl}_2$-action by the
defining property of the Drinfeld coproduct.

This equivariance is the key structural constraint: all
three levels of the $\Ethree$ operations on
$Z^{\mathrm{der}}_{\mathrm{ch}}(V_{\mathrm{EK}})$ are
$\mathfrak{sl}_2$-equivariant maps between representations.
The equivariance arguments of
Proposition~\ref{prop:e3-explicit-sl2} therefore apply
verbatim to the quantum case.

\begin{proposition}[{$\Ethree$ operations on
$Z^{\mathrm{der}}_{\mathrm{ch}}(V_{\mathrm{EK}})$}]
\label{prop:e3-ek-quantum}
Let $V_{\mathrm{EK}}$ be the Etingof--Kazhdan quantum
vertex algebra for $\mathfrak{sl}_2$ at KZ coupling
$h_{\mathrm{KZ}} = 1/(k+2) \neq 0$
\textup{(}Example~\textup{\ref{ex:ek-qvoa}}\textup{)}.
The $\Ethree$-algebra structure on the derived chiral
centre $Z^{\mathrm{der}}_{\mathrm{ch}}(V_{\mathrm{EK}})$
coincides with the classical $\Ethree$ structure of
Proposition~\textup{\ref{prop:e3-explicit-sl2}}: the
cup product and Gerstenhaber bracket vanish, and the
$\Pthree$ bracket is
\begin{equation}\label{eq:ek-p3-bracket}
  \{X, Y\}_q \;=\; h_{\mathrm{KZ}} \cdot (X, Y)
  \qquad
  \textup{for } X, Y \in \HH^1 = \mathfrak{sl}_2,
\end{equation}
with all other brackets zero. In particular, the
vertex $R$-matrix $S(z)$ does not renormalise the
$\Pthree$ bracket on cohomology.
\end{proposition}

\begin{proof}
The argument has four ingredients: equivariance,
one-dimensionality of the obstruction space,
exactness of the rational classical $r$-matrix, and a
$\GRT_1$ consistency check.

\textit{Step 1: cup product and BV operator.}
The cup product $\mu \colon \HH^p \otimes \HH^q
\to \HH^{p+q}$ is a graded-commutative product on
$\mathfrak{sl}_2$-modules over $\CC[[h_{\mathrm{KZ}}]]$.
For $X, Y \in \HH^1
= \mathfrak{sl}_2[[h_{\mathrm{KZ}}]]$,
the product $\mu(X, Y) \in \HH^2
= \CC[[h_{\mathrm{KZ}}]]$
is an $\mathfrak{sl}_2$-equivariant antisymmetric
bilinear form on the adjoint representation:
$\Lambda^2(\mathrm{ad})^{\mathfrak{sl}_2} = 0$,
so $\mu(X, Y) = 0$ over $\CC[[h_{\mathrm{KZ}}]]$.

The BV operator $\Delta \colon \HH^n \to \HH^{n-1}$
is $\mathfrak{sl}_2$-equivariant. The components
$\Delta|_{\HH^1} \colon \mathrm{ad} \to \CC$ and
$\Delta|_{\HH^2} \colon \CC \to \mathrm{ad}$ vanish
by Schur's lemma. Hence $\Delta = 0$ on the derived
centre, and the BV relation gives
$[X, Y] = 0$ for all $X, Y \in \HH^1$, just as in the
classical case.

\textit{Step 2: the $\Pthree$ bracket is proportional
to the Killing form.}
For $X, Y \in \HH^1$, the bracket
$\{X, Y\}_q \in \HH^0 = \CC[[h_{\mathrm{KZ}}]]$ is an
$\mathfrak{sl}_2$-equivariant symmetric bilinear
form on the adjoint representation:
$\Sym^2(\mathrm{ad})^{\mathfrak{sl}_2} \cong \CC$.
